% Secondary workload sensitivity on unchanged frozen Q. Numeric macros are
% generated directly from the sealed readout; this file does not change it.
% Generated by experiments/export_thesis_evaluation_20261006.py; do not hand-edit.
\newcommand{\EvalDynamicRows}{%
Geo-I L20, hiện tại & 87,85 & 85,92 & 91,35 & 785,42 \\
Geo-I L20, cache 60 s & 88,37 & 86,44 & 91,82 & 785,42 \\
Geo-I L30, hiện tại & 91,44 & 89,73 & 94,24 & 1023,95 \\
Geo-I L30, cache 60 s & 91,78 & 90,12 & 94,57 & 1023,95 \\
GPS thật L20 & 99,26 & 99,23 & 99,40 & 156,84 \\
GPS thật L30 & 99,73 & 99,71 & 99,81 & 204,55 \\
Catalogue cũ, loại trạng thái hết hạn & 1,21 & 9,78 & 0,00 & 4,01 \\
Catalogue đầy đủ, cập nhật hiện tại & 100,00 & 100,00 & 100,00 & 126,89 \\
\textbf{Catalogue cũ, không hợp lệ*} & 50,49 & 58,88 & 47,95 & 4,01 \\
}

\section{Kiểm tra trạng thái dịch vụ thay đổi theo thời gian}
\label{sec:current-dynamic-status}

Catalogue tĩnh nhỏ cho phép tải toàn bộ POI rồi tìm kiếm tại thiết bị. Để kiểm
tra một điều kiện gần với dịch vụ có trạng thái thay đổi, thí nghiệm bổ sung
trạng thái \emph{đang khả dụng} cho mỗi POI theo từng epoch công khai 60 giây.
Nó trả lời câu hỏi hẹp: khi phản hồi bị giới hạn ở top-L, thiết bị có giữ
được chất lượng tìm kiếm và tránh dùng trạng thái hết hạn hay không? Đây là
kiểm tra độ nhạy trên dữ liệu đã xem, không phải một xác nhận mới hoặc phép
đo riêng tư bổ sung.

\subsection{Hợp đồng phản hồi và cache còn hiệu lực}

Dịch vụ lấy top-L POI \textbf{theo vị trí tĩnh trong từng category}, rồi gắn
bit khả dụng hiện tại cho các ID đó. Nó không lọc POI khả dụng trước khi lấy
top-L. Thiết bị vẫn gửi tất cả category theo lịch cố định; purpose, bán kính
và đích riêng chỉ dùng để xếp hạng tại thiết bị. L20 và L30 dùng đúng Q, lịch,
neo và ledger đã lưu; không tạo thêm điểm bảo vệ hoặc lần đọc GPS. Trường
epoch và yêu cầu trạng thái được thêm công khai vào bản tin dịch vụ.

Hai chính sách local được so sánh. \textbf{Current-only} dùng các POI nhận ở
event hiện tại. \textbf{Cache epoch60} cộng dồn các phản hồi đã nhận trong
cùng epoch tuyệt đối, nên không thêm request hoặc byte. ID và thông tin tĩnh
có thể được giữ lâu hơn, nhưng bit trạng thái hết hạn khi qua epoch. Epoch
được tính bằng thời gian tuyệt đối của tám phiên, không bắt đầu lại từ 0
ở mỗi chuyến. Thiết bị chỉ biết bit đã nhận qua API.

Cần phân biệt ba trường hợp: POI chưa được truy hồi; có thông tin tĩnh nhưng
\emph{chưa biết trạng thái hiện tại}; và được biết là không khả dụng. Không
suy ``không khả dụng'' từ việc POI vắng trong phản hồi. Câu trả lời live chỉ
chứa POI có trạng thái hiện tại đã biết và đang khả dụng. Tham chiếu là top-5
POI hiện đang khả dụng theo từng purpose/category. Tham chiếu rỗng giữ N/A;
nếu tham chiếu tồn tại nhưng không có ứng viên hợp lệ, Recall bằng 0.

\subsection{Workload và kết quả có kiểm soát}

Workload được khóa trước khi chấm: 418 POI, 88 epoch quan sát và xác suất
khả dụng 0,5, sinh bằng hash của ID, epoch và seed cố định. Replay gồm toàn
bộ 24 family test đã xem, mỗi family ba draw và tám phiên, tổng cộng 18.114
event cho mỗi arm. Tất cả family/draw dùng cùng một thế giới trạng thái.
Recall được gộp theo category có tham chiếu, event, draw và family; cuối
cùng lấy trung bình đều bốn purpose, như cách gộp utility đã định nghĩa.

\begin{table}[htbp]
\centering\small
\caption{Kiểm tra trạng thái tổng hợp trên cùng Q. Recall là \%; cold là
phiên đầu, tail là thời gian tương đối $t\ge400$ giây. JSON MB dùng
$10^6$ byte, tính cả metadata và trạng thái.}
\label{tab:current-dynamic-status}
\begin{tabularx}{\textwidth}{Yrrrr}
\toprule Arm & Recall & Cold & Tail & JSON MB \\\midrule
\EvalDynamicRows
\bottomrule
\end{tabularx}
\end{table}

L30 current-only đạt \EvalDynamicLThirtyCurrent{}\%, tăng
\EvalDynamicGain{} điểm phần trăm so với L20, với tổng JSON tăng
\EvalDynamicCostGrowth{}\%. Cache epoch60 đưa L30 lên
\EvalDynamicLThirtyCached{}\% mà không tăng truy vấn hoặc byte. Đây là
utility dưới workload mới; L30 vẫn là độ sâu đã khóa từ phép thử trước,
không được chọn lại bằng kết quả trạng thái này. Purpose trong bán kính có
\EvalDynamicRadiusDefined{}/18.114 event được định nghĩa; các event không
có tham chiếu không được cộng điểm hoàn hảo.

Mọi arm vận hành hợp lệ đều trả về 0 POI không khả dụng và 0 POI có trạng
thái hiện tại chưa biết. Control catalogue cũ \emph{safe} có Recall thấp vì
từ chối bit hết hạn. Control \emph{unsafe} cố ý dùng lại bit cũ và được đánh
dấu \textbf{không hợp lệ}: trong 1.279.509 item trả về theo bốn purpose,
623.325 hiện không khả dụng và 1.264.314 chưa có trạng thái hiện tại đã biết.
Hai số này có phần giao nhau; không xem Recall của arm đó là chất lượng câu
trả lời live an toàn.

Control tải toàn catalogue với bit hiện tại đạt 100\%, không cần gửi vị trí
và dùng ít JSON hơn L30. Nó có API mạnh hơn provider top-L, nhưng phải được
giữ trong so sánh nếu ứng dụng cho phép bulk. Vì vậy trạng thái động chưa
giải quyết khoảng trống catalogue nhỏ/tải toàn bộ; muốn lập luận dịch vụ từ
xa là cần thiết còn phải có hợp đồng provider và workload thực phù hợp.

\subsection{Kiểm chứng và giới hạn diễn giải}

Checker độc lập đã PASS cho 72 block và 163.026 dòng arm/event: tái dựng
trạng thái, thứ tự tham chiếu đường có hướng, cache nhận theo thời gian,
partition unknown, mẫu số và chi phí. Bản V1 cùng khai báo trước scoring và
failure do thứ tự serialization được giữ nguyên. V2 sửa cách ghép bằng
khóa event/arm, được khai báo rõ sau scoring; không đổi nguồn runner, thế
giới trạng thái, rows hay readout. Hồ sơ ở
\path{artifacts/benchmarks/dynamic_provider_status_20261006_v1/}.

Seed và fixture đầy đủ công khai trong artifact nghiên cứu để tái lập,
nhưng bị loại khỏi API client của thí nghiệm. Client được cấp fixture có
thể tự tính mọi trạng thái. Do đó phép thử kiểm tra freshness và luồng dữ
liệu dưới giả định API, không chứng minh trạng thái thực không thể dự đoán
hoặc phải gọi server. Một thế giới tổng hợp không đại diện nhiều quá trình
khả dụng độc lập. Byte chỉ là compact JSON ứng dụng; chưa đo mạng, latency,
pin, traffic động hoặc giá dịch vụ. Kết quả này giữ nguyên nền tảng Geo-I,
bổ sung điều kiện dùng cache live, và không thay thế bằng chứng xác nhận
utility tĩnh hay chứng nhận riêng tư mới.
